\chapter*{Abstract}
Retrieval Augmented Generation significantly improves Large Language Models but faces issues involving complicated syntactic ambiguities and loss of data when creating knowledge graphs. To tackle these problems, the Quantum Retrieval Augmented Generation is suggested in the context of Farm Fetching architecture by combining the efficient classical pipeline designed for linear processing with the quantum processor responsible for non-linear structure analysis. First, a classical approach evaluates 130,000 documents in terms of MongoDB vector search, dynamic Neo4j Knowledge Graphs, and baseline language model. This part of the research involves the introduction of Constructive Hallucination, a technique designed to validate and include artificial objects in order to enhance dynamic knowledge graphs. While various tests, including Thematic Correlator, showed superior performance of classical models in terms of semantics over the quantum systems, the classical parser cannot deal with deep structural ambiguity because of local attachment bias. For these difficult situations, a Quantum NLP Module performs an Ambiguity Controller, which is responsible for the redirection of requests. Using grammatical dependencies as input, quantum parameters are built into the quantum circuit, which uses entanglement to process the non-linear structures. The accuracy gain achieved in this way was 46 percent compared to the classical heuristic solution. Therefore, the Utility Threshold was identified, meaning that the combination of classical efficiency and quantum-nativeness provides a solid foundation for future studies' evaluation systems.